\section{Labelled expanders and their finite cycle spaces}\label{sec:data}

Fix a positive \(\lambda<1/24\) so small that
\begin{equation}\label{eq:lambda-choice}
12(L_0+1)\lambda\le\frac1{16},
\end{equation}
where \(L_0\) is the universal constant from
\eqref{eq:root-constants}. Put \(\eta=1/100\).

\begin{proposition}[Choice of data]\label{prop:data}
There are labelled graphs \((\Theta_i)_{i\ge1}\) satisfying all the
hypotheses of Section~\ref{sec:geometry}, with the following additional
properties. Write \(n_i=|\operatorname{Vert}(\Theta_i)|\),
\(k_i=\lfloor n_i/2\rfloor\), and let \(\sigma_i\) be the girth.
There are constants \(\rho,\kappa>0\), independent of \(i\), such that:
\begin{enumerate}[label=\textup{(\roman*)}]
\item On the orthogonal complement of constants in
\(\ell^2(\operatorname{Vert}(\Theta_i))\),
\(dI-\Adj_{\Theta_i}\ge\rho I\), and for every vertex subset \(W\),
\begin{equation}\label{eq:expansion}
|\partial W|\ge\kappa\min\{|W|,n_i-|W|\}.
\end{equation}
Here \(\partial W\) is the unoriented edge boundary in \(\Theta_i\).
\item With
\begin{equation}\label{eq:c-definition}
c=\frac\eta2\min\{\kappa,1\},
\end{equation}
we have \(n_i>100\), \(\sigma_i>\max\{16,20/c\}\),
\(\lambda\sigma_i>3\), and
\(1<\lfloor\lambda\sigma_i\rfloor<\lfloor\lambda\sigma_{i+1}\rfloor\).
\item Every induced intrinsic ball of radius at most \(\sigma_i/16\)
is a tree with fewer than \(\eta n_i\) vertices.
\item The integers \(k_i\) strictly increase and
\(k_{i+1}-k_i\to\infty\).
\end{enumerate}
Consequently every ball \(B_P\) in Lemma~\ref{lem:overlap}, for every
permitted root and chosen root geodesic, is an induced tree satisfying
\begin{equation}\label{eq:small-ball}
r_P\le\sigma(P)/16,\qquad |B_P|<\eta|P|.
\end{equation}
\end{proposition}

\begin{proof}
Start with a fixed-degree family of connected regular expanders of
unbounded order and girth at least a positive constant times logarithmic
order, as supplied by the Lubotzky--Phillips--Sarnak construction
\cite[Theorems~3.4 and~4.1]{LPS1988}; see also \cite[Section~2.4]{Osajda2020}. Discarding small
graphs leaves simple graphs. Their Laplacians have a common gap
\(\rho>0\) off the constants. Bipartiteness, if present, causes no
problem: the adjacency eigenvalue \(-d\) gives Laplacian eigenvalue
\(2d\), not another zero eigenvalue.

For \(w=|W|\), the centered indicator of \(W\) has squared norm
\(w(n-w)/n\) and Laplacian energy \(|\partial W|\). Thus
\[
|\partial W|\ge\rho\frac{w(n-w)}n
\ge\frac\rho2\min\{w,n-w\},
\]
so we may choose \(\kappa=\rho/2\). A ball containing at most half the
vertices has an outside vertex boundary of size at least \(\kappa/d\)
times its size. Balls grow by a fixed factor until they contain more than
half the graph. Any two such balls intersect, giving a uniform
\(O(\log n)\) diameter bound. Hence the diameter/girth ratio is bounded.

An induced ball of integer radius \(r\) with \(2r+1<\sigma\) is a
tree: an edge outside a shortest-path tree would create a cycle of
length at most \(2r+1\). Its order at \(r\le\sigma/16\) is at most
\[
1+d\sum_{j=0}^{\lfloor\sigma/16\rfloor-1}(d-1)^j.
\]
Regularity and the same tree property at radius
\(\lfloor(\sigma-2)/2\rfloor\) give the lower bound
\[
n\ge1+d\sum_{j=0}^{\lfloor(\sigma-2)/2\rfloor-1}(d-1)^j.
\]
The ratio tends to zero with \(\sigma\). A subsequence therefore satisfies
(ii)--(iv), including (iii); growing order allows the gaps in (iv) to be
chosen arbitrarily large.

The resulting sequence has bounded degree, growing girth, bounded
diameter/girth ratio, and the required strictly increasing floors.
Osajda's labelling theorem \cite[Theorem~2.7]{Osajda2020}, applied to the
already chosen \(\lambda\), gives a reduced labelling over one finite
alphabet with no repeated long path word. Its conclusion for
nonbacktracking paths includes the simple-path condition used here.
The associated graphical presentation has isometrically embedded
relators by \cite[Lemma~3.1 and Theorem~3.2]{Osajda2020}.
This supplies the data of Section~\ref{sec:geometry}. Finally,
\eqref{eq:lambda-choice}, \eqref{eq:overlap-radius}, and (iii) give
\eqref{eq:small-ball}.
\end{proof}

\paragraph{Order of choices.}
The constants used above are chosen in the order recorded in
Table~\ref{tab:parameters}. In particular, the geometric constant is
universal before the small-cancellation parameter is selected. Neither
the alphabet nor the code field is fixed while choosing the expanders.

\begin{table}[ht]
\centering
\small
\begin{tabularx}{\linewidth}{@{}lX@{}}
\toprule
Choice & Requirement and subsequent use \\
\midrule
$D_0,t,L_0$ & Universal thinness and root-path constants from
Proposition~\ref{prop:cone-geometry} and \eqref{eq:root-constants}.\\
$\lambda$ & $0<\lambda<1/24$ and
$12(L_0+1)\lambda\le1/16$, so every overlap ball has radius at most
one sixteenth of the girth.\\
$d,\rho,\kappa$ & One fixed-degree expander family; take
$\kappa=\rho/2$. Set $\eta=1/100$ and
$c=(\eta/2)\min\{\kappa,1\}$.\\
$(\Theta_i)$ & Pass to a subsequence with
$\sigma_i>\max\{16,20/c\}$, small tree balls, strictly increasing
$\lfloor\lambda\sigma_i\rfloor$, and diverging gaps between
$k_i=\lfloor n_i/2\rfloor$.\\
$S$ & A single finite alphabet supplied by the labelling theorem for
this sequence and this value of $\lambda$.\\
$p,C_i$ & Choose a prime $p>2$ encoding all $d$-element sets of outgoing labels;
then choose the $k_i$-dimensional codes in Lemma~\ref{lem:codes}.\\
$\varepsilon$ & The uniform gap is
$\min\{\rho,c^2/(8d)\}$, as proved in
Lemma~\ref{lem:spectral-gap}.\\
\bottomrule
\end{tabularx}
\caption{Parameter choices for $\Gsur$. Each row uses only earlier
choices; the last row records the resulting spectral estimate.}
\label{tab:parameters}
\end{table}

From now on all these data, the group \(\Gamma\), its Cayley graph
\(X\), and the coned graph \(\widehat X\) are fixed. By
Proposition~\ref{prop:cone-geometry}, \(X\) is simple. We identify a copy
with its vertex set when discussing supports and specify its edges when
discussing chains.

For a graph \(Y\), let \(Z_1(Y;\F_2)\) be its ordinary
\emph{finite-support} cycle space: its elements are finite sums of edges
with coefficients in \(\F_2\), with even incidence at every vertex.
Thus no completed chain spaces are used in the next statement.

\begin{lemma}[Direct sum of copy cycles]\label{lem:cycle-sum}
The inclusions of the copies give an internal algebraic direct sum
\begin{equation}\label{eq:cycle-sum}
Z_1(X;\F_2)=\bigoplus_P Z_1(P;\F_2).
\end{equation}
\end{lemma}

\begin{proof}
Every finite cycle is a sum of closed walks. The label of any such walk
is trivial in \(\Gamma\), so a finite van Kampen diagram over the
presentation expresses its
edge chain, modulo two, as the sum of the chains of its bounded-face
walks. Each lies in a copy. This proves spanning.

Suppose finitely many nonzero cycles in distinct copies sum to zero.
Choose a group-vertex root \(R\) and a participating copy \(P\) whose
center is farthest from \(R\). By Lemma~\ref{lem:overlap}, all its
intersections with the other participating copies lie in one ball
\(B_P\), which is an induced tree by \eqref{eq:small-ball}. Every
nonzero edge coefficient in its cycle must cancel against an edge of
another participating copy. Both endpoints therefore lie in \(B_P\),
and inducedness puts the edge itself in that tree. A tree has no nonzero
finite cycle, a contradiction.
\end{proof}
